%% 09_conclusion.tex
%% Section 9: Conclusion

\section{Conclusion}
\label{sec:conclusion}

This paper develops the first integrated framework for assessing how coal stranded asset risk propagates through the banking sector to the macroeconomy in Indonesia, the world's largest thermal coal exporter. Using firm-level data from 35 listed coal companies and 38 listed banks, we construct a three-model framework that links bottom-up stranded asset valuations to bank-level credit losses and aggregate macroeconomic impacts.

Our key findings are as follows. First, under a 2\textdegree C-aligned climate transition, aggregate coal stranded values reach \$79.7 billion (90\% CI: \$72.7--88.0 billion from Monte Carlo simulation with 10,000 draws), rising to \$128.3 billion (90\% CI: \$117.1--141.0 billion) under a 1.5\textdegree C pathway. Stranding risk is highly concentrated: the top five companies account for over 75\% of total stranded value. Second, the banking sector stress test reveals total credit losses of \$1.8 billion under the 2\textdegree C scenario and \$3.1 billion under 1.5\textdegree C, with losses concentrated in a small number of large banks. While no bank breaches regulatory capital thresholds, the most exposed D-SIB (BMRI) absorbs losses equivalent to 42\% of annual earnings, highlighting the idiosyncratic vulnerability created by concentrated coal lending. Third, the macro-financial transmission framework estimates total GDP impacts of 0.32\% under the 2\textdegree C pathway and 0.54\% under 1.5\textdegree C, with the credit channel (bank deleveraging) accounting for approximately 60\% of the total effect.

These findings contribute to the literature in three ways. To the stranded assets literature, we contribute the most granular firm-level valuation of coal stranding in any single country, with Monte Carlo uncertainty quantification that moves beyond deterministic point estimates. To the climate stress testing literature, we contribute the first bottom-up stress test of the Indonesian banking system using actual bank-coal credit relationships rather than sectoral aggregates. To the macro-financial literature, we contribute a calibrated three-channel transmission framework that decomposes the aggregate GDP impact of coal stranding into fiscal, credit, and market components.

The policy implications are clear. Indonesia's regulators, OJK and Bank Indonesia, should integrate coal transition risk into supervisory stress testing, enhance fossil fuel exposure disclosure requirements, and consider climate-adjusted capital requirements for coal-exposed banks. The JETP's success depends on managing the financial vulnerabilities we document: an orderly, well-signalled transition that allows banks and firms to adjust gradually will minimise the credit channel amplification that our analysis identifies as the dominant transmission mechanism. Conversely, a disorderly transition risks triggering the bank deleveraging spiral that converts firm-level stranding into a macroeconomic headwind.

More broadly, the coal-banking-sovereign nexus we document is likely replicated in other coal-dependent emerging economies, including South Africa, India, Colombia, and the Philippines. Our framework provides a replicable template for country-specific assessments, and we encourage future research to extend the analysis to these contexts. Additional directions for future work include incorporating dynamic transition pathways (rather than static scenario comparisons), modelling second-round effects through interbank contagion and sovereign-bank feedback loops, and extending the analysis to the full coal value chain including power generation, logistics, and services.

The urgency of this research agenda is underscored by the pace of the global energy transition. As renewable energy costs continue to fall and climate policy ambition rises, the question for coal-dependent economies is not whether stranding will occur, but when and how fast. Our analysis demonstrates that the financial and macroeconomic consequences of coal stranding in Indonesia are material but manageable, provided that policymakers act early to build resilience in the banking sector, diversify fiscal revenues, and ensure that the transition is orderly, equitable, and well-financed.
